\begin{theorem}[零温冻结与 \texorpdfstring{$\infty$}{infty} 预算端点塌缩的不交换性]\label{thm:conclusion-freezing-infinity-budget-noncommutation}
\leanverified{paper\_conclusion\_freezing\_infinity\_budget\_noncommutation}
沿用小节 \ref{subsec:conclusion-maxfiber-freezing-capacity-entropy-separation} 的记号，并假设严格顶端间隙
\[
\alpha_\ast^{(2)}<\alpha_\ast
\]
成立。则对每个 \(a>a_{\mathrm c}\)，冻结相同时保留顶端纤维的指数高度与基态简并熵，即
\[
P_a=a\alpha_\ast+g_\ast,
\qquad
\lim_{m\to\infty}\frac{1}{m}H(\pi_{m,a})=g_\ast,
\]
且 escort 测度 \(\pi_{m,a}\) 在最大纤维集合
\[
A_m^{\max}:=\{x\in X_m:\ d_m(x)=M_m\}
\]
上指数集中。另一方面，若把连续容量曲线在预算参数上补入端点 \(T=\infty\)，则
\[
\mathcal C_m^{\mathrm{cont}}(\infty)=\sum_{x\in X_m}d_m(x)=2^m,
\qquad
R_m(\infty)=0,
\]
从而该单点只保留总质量，而不再保留 \((M_m,K_m,M_m^{(2)})\) 的任何谱几何信息。

因此，零温冻结极限保留基态简并熵 \(g_\ast\) 及最大纤维几何，而 \(\infty\)-预算端点却把这些信息全部抹去；这两种奇异极限在结构上并不交换。
\end{theorem}
\begin{proof}
由定理 \ref{thm:conclusion-frozen-escort-shannon-renyi-collapse}，对每个 \(a>a_{\mathrm c}\) 都有
\[
P_a=a\alpha_\ast+g_\ast,
\qquad
\pi_{m,a}(A_m^{\max})\ge 1-\exp\!\bigl(-m\Delta(a)+o(m)\bigr),
\]
并且
\[
\lim_{m\to\infty}\frac{1}{m}H(\pi_{m,a})=g_\ast.
\]
这表明冻结相不仅保留顶端纤维的指数高度 \(\alpha_\ast\)，而且保留其指数简并度 \(g_\ast\)。

另一方面，由定理 \ref{thm:conclusion-capacity-deficit-mellin-bernstein-completion}，
\[
R_m(T)=2^m-\mathcal C_m^{\mathrm{cont}}(T).
\]
对充分大的 \(T\) 有 \(R_m(T)=0\)，故把预算参数补入端点 \(T=\infty\) 时必有
\[
\mathcal C_m^{\mathrm{cont}}(\infty)=2^m,
\qquad
R_m(\infty)=0.
\]
而定理 \ref{thm:conclusion-capacity-deficit-mellin-bernstein-completion} 同时说明：真正携带完整多重度谱几何的是整个有限预算族 \(\{\mathcal C_m^{\mathrm{cont}}(T)\}_{T>0}\)，不是单点 \(T=\infty\)。因此，\(\infty\)-预算端点只能记住总质量 \(2^m\)，不能恢复 \(M_m\)、\(K_m\)、\(M_m^{(2)}\)，也就不能恢复 \(g_\ast\)。

于是，一侧的零温冻结极限保留顶端几何与简并熵，另一侧的 \(\infty\)-预算端点却把它们全部抹平；两种极限在结构上并不交换。证毕。
\end{proof}

\begin{conjecture}[冻结阈值与 Hankel 共振阈值的一致性]\label{conj:conclusion-freezing-hankel-threshold-coincidence}
沿用定理 \ref{thm:fold-pressure-freezing-threshold} 的冻结阈值记号，并记
\[
q_{\mathrm{fr}}
:=
\min\{q\in\ZZ_{\ge 2}:\ P(q)=q\alpha_\ast+g_\ast\}.
\]
另一方面，沿用正文关于折叠碰撞矩最小整系数递推
\[
S_q(m)=\sum_{i=1}^{d_q}c_iS_q(m-i)
\]
的记号，定义 Hankel 共振阈值
\[
q_{\mathrm H}
:=
\min\{q\in\ZZ_{\ge 2}:\ d_q<2\lfloor q/2\rfloor+1\}.
\]
则在 Zeckendorf 规范核中应有
\[
q_{\mathrm{fr}}=q_{\mathrm H}.
\]
亦即，热力学冻结阈值与最小线性实现的第一次秩塌缩并非两类独立现象，而是同一相变在压力函子与 Hankel 函子下的两种投影。
\end{conjecture}
